\subsection{CAUCHY'S CRITERION}

Let \((x_{t})_{t\in I}\) be a summable family in \(G\) . Then for each neighbourhood \(V\) of the origin in \(G\) there is a finite subset \(J_{0}\) of \(I\) such that,

for each finite subset K of I which does not meet \(J_{0}\) , we have \(s_{K} \in V\) . For \(J = J_{0} \cup K\) is an arbitrary finite subset containing \(J_{0}\) ; let \(s = \sum x_{i}\) and let W be a symmetric neighbourhood of o such that \(W + W \subset V\) ; then, by Definition 1, we may choose \(J_{0}\) such that \(s_{J} \in s + W\) and \(s_{J_{0}} \in s + W\) , which means that \(s_{K} = s_{J} - s_{J_{0}} \in W + W \subset V\) .

Conversely, suppose that the family \((x_{i})\) has this property. Then the image of the filter \(\Phi\) , under the mapping \(J \to s_{J}\) , is a Cauchy filter base in G. For if J is a finite subset containing \(J_{0}\) , and if K denotes \(J \cap \complement J_{0}\) , then \(K \cap J_{0} = \emptyset\) and \(s_{K} = s_{J} - s_{J_{0}}\) , so that \(s_{J} \in s_{J_{0}} + V\) . If \(J'\) is another finite subset containing \(J_{0}\) , then \(s_{J} - s_{J'} \in V + V\) , and the result follows. Consequently:

THEOREM 1 (Cauchy's criterion). In a Hausdorff commutative group G, in order that a family \((x_{t})_{t\in I}\) should be summable it is necessary that, for each neighbourhood V of the origin in G, there is a finite subset \(J_0\) of I such that \(\sum_{i\in K}x_i\in V\) for all finite subsets K of I which do not meet \(J_0\) . This necessary condition is also sufficient if G is complete.

Thus by taking away a (sufficiently large) finite number of terms from the family \((x_{i})\) , every finite partial sum of the surviving subfamily has to be as close to o as we please.

An immediate consequence of the first part of Theorem 1 is the following proposition:

PROPOSITION 1. If the family \((x_{t})\) is summable, then every neighbourhood of o contains all the \(x_{t}\) except for a finite subfamily (in other words, if I is infinite, we have \(\lim x_{t} = 0\) with respect to the filter of complements of finite subsets of I).

\pandocbounded{\includegraphics[keepaspectratio]{images/part002_cfff7d82.jpg}}

This necessary condition for a family \((x_{i})\) to be summable is by no means sufficient in general, even if \(G\) is complete; we shall see many examples later (see Chapter IV, § 7).

COROLLARY. Let \((x_{t})_{t\in I}\) be a summable family in a commutative group whose identity element has a countable fundamental system of neighbourhoods. Then the set of indices \(\iota\) such that \(x_{\iota}\neq 0\) is countable.

Let \((\mathbf{V}_{n})\) be a countable fundamental system of neighbourhoods of 0. If \(\mathbf{H}_{n}\) is the set of all indices i such that \(x_{i} \notin \mathbf{V}_{n}\) , then the set H of indices i such that \(x_{i} \neq 0\) is the union of the sets \(\mathbf{H}_{n}\) , and each of the \(\mathbf{H}_{n}\) is finite by Proposition 1.

This corollary is no longer necessarily valid if we do not assume that the origin has a countable fundamental system of neighbourhoods. * Consider, for example, the product group \(\mathbf{R}^{\mathbb{R}}\) {[}the additive group of all finite real-valued functions of a real variable, with the topology of pointwise convergence (cf.~Chapter X, § 1, no. 3){]}, and let \(f_{a}\) be the element of \(\mathbf{R}\) such that \(f_{a}(a)=1\) and \(f_{a}(x)=0\) if \(x\neq a\) ; then the family \((f_{a})_{a\in\mathbb{R}}\) is summable and its sum is the function whose value at each point of \(\mathbf{R}\) is 1. *

Remark. When G is written multiplicatively, Cauchy's criterion takes the following form: for the family \((x_{i})_{i\in I}\) to be multipliable it is necessary that, for each neighbourhood V of the identity element, there should exist a finite subset \(J_{0}\) of I such that, for each finite subset K of I which does not meet \(J_{0}\) , we have \(\prod_{i\in K}x_{i}\in V\) ; and this condition is sufficient provided that G is complete. We deduce that if I is infinite and \((x_{i})\) is multipliable, then \(\lim x_{i}=1\) with respect to the filter of complements of finite subsets of I; if in addition the identity element has a countable fundamental system of neighbourhoods, then the set of indices i such that \(x_{i}\neq1\) is countable.

